\section{The actual full canonical maximum at the initial NC history}
\label{sec:hironaka-normal-crossings-initial-full-maximum}

Assume the ambient complex dimension $n$ is positive. For an actual NC
carrier $H$, let $M_H=(X,X,\varnothing,I_H,1)$ be its unchanged original
marked datum, and let $V_H(x)$ be the whole value of the actual canonical
family of depth $n+1$ at the reflexive tower. Every assertion in this
section concerns this initial empty boundary and empty history. No
later chronological or companion history is identified with branch count.

\begin{lemma}[The genuine coordinate NC model]
\label{hir:general-nc-standard-coordinate-carrier}
\leanok
\lean{Hironaka.HironakaNormalCrossingsInitialFullMaximum.coordinateCarrier_isNormalCrossings}
For a finite set $S$ of actual coordinate indices, the whole union
$H_S=\bigcup_{i\in S}\{z_i=0\}$ in the lifted complex coordinate model
is a genuine NC carrier.
\end{lemma}
\begin{proof}
Every coordinate hyperplane is closed, so their finite union is closed.
At a point $w$, let $S_w$ be those selected
indices for which $w_i=0$. Restrict the translation chart $z\mapsto z-w$
to the open neighbourhood on which all selected coordinates nonzero
at $w$ remain nonzero. On that actual domain, membership in the whole
original union is equivalent to vanishing of at least one selected
coordinate in $S_w$, and the corresponding translated coordinate has
the same value because $w_i=0$. This is exactly the centred NC
presentation of the whole carrier. An empty selected set gives an
empty union and the empty local branch presentation.
\end{proof}

\begin{lemma}[Whole initial values under an actual NC flattening]
\label{hir:general-nc-initial-full-value-flat-chart}
\leanok
\lean{Hironaka.HironakaNormalCrossingsInitialFullMaximum.initialFullValue_eq_coordinate_model}
Let an actual centred NC flattening at a point of $H$ identify the whole
carrier on its genuine chart domain with the selected coordinate union
$H_S$. Then $V_H(x)=V_{H_S}(0)$, including every lower entry of the
full reflexive invariant.
\uses{hir:general-nc-standard-coordinate-carrier,hir:general-nc-actual-local-snc-branches}
\end{lemma}
\begin{proof}
The genuine flattening and its genuine inverse give an actual
biholomorphism from the open flat domain to the lifted open target.
It takes the original point to zero and takes the whole carrier trace
to the literal trace of $H_S$. At every point of these domains,
composition with this actual biholomorphism and its inverse identifies
the entire holomorphic vanishing-ideal predicates: a holomorphic germ
vanishes on one trace if and only if its transported germ vanishes on
the other. Consequently the actual pulled-back marked data have the
same full ambient carrier, empty boundary, whole original ideal and
marking one. Their actual invariant domains follow from the proved NC
original-domain theorem and actual open/biholomorphic transport.

The reflexive history is genuinely canonical without any additional
condition. Apply the full invariant's proved open restriction law to
the original flat domain, its whole tower-isomorphism law to the actual
flattening biholomorphism, and its open restriction law to the coordinate
target. The original point belongs to the original marked cosupport
and zero belongs to the coordinate marked cosupport. These three exact
value equalities, with the whole marked-data identities just proved,
give the claimed equality of full values.
\end{proof}

\begin{lemma}[Full initial values agree on the genuine highest branch stratum]
\label{hir:general-nc-initial-full-highest-stratum-constant}
\leanok
\lean{Hironaka.HironakaNormalCrossingsInitialFullMaximum.initialFullValue_eq_of_highest_branch_count}
Suppose every original branch count is at most $r>0$. Any two points of
branch count $r$ have the same whole full initial canonical value.
\uses{hir:general-nc-initial-full-value-flat-chart,hir:general-nc-highest-stratum-branch-adapted-coordinates}
\end{lemma}
\begin{proof}
The actual highest-stratum coordinate producer gives a centred
flattening at each point in which the original carrier is the union of
the last $r$ coordinate hyperplanes. Both actual charts therefore use
the same selected coordinate set, without choosing or comparing global
branch labels. Positivity of $r$ puts both points in the original
carrier. The preceding whole-value transport sends each full initial
value to the same coordinate-model value at zero. Transitivity proves
their equality, with all lower entries retained.
\end{proof}

\begin{theorem}[The actual initial full maximum is the whole highest branch stratum]
\label{hir:general-nc-initial-whole-full-maximum}
\leanok
\lean{Hironaka.HironakaNormalCrossingsInitialFullMaximum.initialFullMaximum_eq_highest_stratum}
Suppose the positive bound $r$ is attained and bounds every original
branch count. The actual whole maximum locus of $V_H$ on its actual
original cosupport is exactly $\{x:r\leq\nu_H(x)\}$.
\uses{hir:general-nc-initial-full-highest-stratum-constant,hir:general-nc-actual-prefix-first-entry}
\end{theorem}
\begin{proof}
The actual original cosupport is exactly $H$. Its genuine first full
entry is the intrinsic original branch count. Fix an attained point
of count $r$. Every other point of count $r$ has its same whole full
value by the preceding lemma. At every remaining point of $H$, the
first entry is strictly smaller than $r$, so the actual lexicographic
full value is strictly smaller, regardless of its later entries.
Thus every point of the whole $r$-branch stratum is a full maximum.
Conversely a full maximum cannot have smaller branch count, since its
value would be strictly below the attained reference value. All counts
are bounded by $r$, so the maximum lies in precisely that stratum.
Both set inclusions retain the entire original cosupport.
\end{proof}

\begin{theorem}[An actual first full canonical blowup from only the original NC carrier]
\label{hir:general-nc-original-first-full-canonical-blowup}
\leanok
\lean{Hironaka.HironakaNormalCrossingsInitialFullMaximum.exists_initial_full_canonical_blowup}
Every nonempty actual original NC carrier internally supplies an attained
positive highest branch count, its nonempty whole highest stratum and
a genuine blowup of that stratum. The actual one-step tower is a test
tower for the unchanged original marked datum and is canonical for its
whole full invariant.
\uses{hir:general-nc-initial-whole-full-maximum,hir:general-nc-highest-stratum-actual-blowup}
\end{theorem}
\begin{proof}
The finite intrinsic count range, bounded by $n$, supplies an attained
positive maximum from nonemptiness of the original carrier. The actual
highest-stratum producer gives its genuine smooth centre and actual
blowup. The preceding theorem identifies that very centre with the
entire actual initial full maximum on the same original cosupport.
The centre is smooth, lies in the marking-one cosupport and has normal
crossings with the empty old boundary, proving the actual first test
step. The actual full canonical-history step constructor uses precisely
the derived full-maximum equality and the reflexive initial history.
No compactness, supplied maximum, value formula, centre, finite history
or termination certificate is an input. At branch count one the
genuine blowup may be Cartier; no differential strict carrier is erased
by assertion. No second or later full canonical step is claimed.
\end{proof}
